\section{Exact residue-fiber orthogonality in the physical packet}
\label{sec:tpc42-residue-orthogonality}

We now apply the abstract decomposition to the actual terminal vectors.
The decisive fact is elementary but rigid: the CRT residue direction and
the physical orbit-coordinate direction are transverse.

\subsection{Fixed-alias Hilbert atoms}

For a retained alias \(k\in\cI_B\), put
\begin{equation}
 A_k(r)=\log(r+h+kM)\Phi_0(r+h+kM),
 \label{eq:tpc42-Akr}
\end{equation}
interpreted as zero outside the positive terminal support, and define
\(\mathbf G_k^L(r)\in\ell^2\{(\gamma,j)\}\) by
\begin{equation}
 [\mathbf G_k^L(r)]_{\gamma,j}
 =\sum_{\alpha:m_\alpha j=r}
   c^L_{\alpha,\gamma,j}A_k(r).
 \label{eq:tpc42-fixed-k-G}
\end{equation}
For fixed \((j,r)\), at most one row occurs: \(r/j\), when integral,
determines \(m_\alpha\), and source separation determines \(\alpha\).
Every sum containing \(\mu(r+h+kM)\) below is restricted to the positive
terminal support, so no value of \(\mu\) at a nonpositive integer is used.
The fixed-alias column is exactly
\begin{equation}
 \mathscr T_k^L
 =-\sum_r\mu(r+h+kM)\mathbf G_k^L(r).
 \label{eq:tpc42-fixed-k-compression}
\end{equation}

\begin{theorem}[Physical residue-fiber orthogonality]
\label{thm:tpc42-fixed-alias-orthogonality}
Fix \(k\in\cI_B\).  If \(r\ne r'\) and
\begin{equation}
 r\equiv r'\pmod M,
 \label{eq:tpc42-same-base-residue}
\end{equation}
then
\begin{equation}
 \boxed{
 \langle\mathbf G_k^L(r),\mathbf G_k^L(r')\rangle=0.}
 \label{eq:tpc42-G-orthogonal}
\end{equation}
The conclusion is uniform in the terminal cutoff, the M\"obius values,
and the actual masks.  It holds identically for the right vectors.
\end{theorem}

\begin{proof}
A nonzero term in the inner product would require a shared coordinate
\((\gamma,j)\).  There would then be rows \(\alpha,\beta\) such that
\begin{equation}
 r=m_\alpha j,
 \qquad
 r'=m_\beta j.
 \label{eq:tpc42-shared-coordinate-products}
\end{equation}
By \eqref{eq:tpc42-same-base-residue},
\begin{equation}
 M\mid(m_\alpha-m_\beta)j.
 \label{eq:tpc42-M-divides-row-time}
\end{equation}
Every \(j\in\cJ\) is positive and smaller than each prime factor of
\(M\), so \((j,M)=1\).  Hence \(M\mid m_\alpha-m_\beta\).  The strict row
gap in \eqref{eq:tpc42-no-wrap-data} forces
\(m_\alpha=m_\beta\), and source separation gives
\(\alpha=\beta\).  Equation \eqref{eq:tpc42-shared-coordinate-products}
then gives \(r=r'\), a contradiction.  Vanishing cutoffs or coefficients
only delete coordinates and cannot invalidate the argument.
\end{proof}

It is convenient to index by the target residue.  For every
\(a\in\Z/M\Z\), define
\begin{equation}
 \mathbf V_{k,a}^L
 :=-\sum_{\substack{r\equiv a-h\!\!\!\pmod M}}
 \mu(r+h+kM)\mathbf G_k^L(r).
 \label{eq:tpc42-fixed-k-residue-packet}
\end{equation}
Define the actual fixed-alias atomic diagonal
\begin{equation}
 \cD_{k,L}
 :=\sum_{\gamma,j,\alpha}
 |c^L_{\alpha,\gamma,j}|^2
 |\mu(n_{\alpha,j,k})|^2
 |A_k(m_\alpha j)|^2.
 \label{eq:tpc42-fixed-k-actual-diagonal}
\end{equation}

\begin{corollary}[Exact packed vector variance]
\label{cor:tpc42-fixed-alias-packets}
For every fixed alias,
\begin{align}
 \mathscr T_k^L
 &=\sum_{a\bmod M}\mathbf V_{k,a}^L,
 \label{eq:tpc42-fixed-k-synthesis}\\
 \sum_{a\bmod M}\|\mathbf V_{k,a}^L\|^2
 &=\cD_{k,L},
 \label{eq:tpc42-fixed-k-packed-diagonal}\\
 \|\mathscr T_k^L\|^2
 &=\cD_{k,L}
   +\sum_{\substack{a,b\bmod M\\a\ne b}}
    \langle\mathbf V_{k,a}^L,\mathbf V_{k,b}^L\rangle.
 \label{eq:tpc42-fixed-k-cross-residue}
\end{align}
The second sum in \eqref{eq:tpc42-fixed-k-cross-residue} is an ordered
sum and is real after pairing \((a,b)\) with \((b,a)\).
\end{corollary}

\begin{proof}
The first identity partitions \eqref{eq:tpc42-fixed-k-compression} by
target residue.  The vectors inside each packet are pairwise orthogonal by
\cref{thm:tpc42-fixed-alias-orthogonality}; hence
\begin{equation}
 \|\mathbf V_{k,a}^L\|^2
 =\sum_{r\equiv a-h\, (M)}
  |\mu(r+h+kM)|^2\|\mathbf G_k^L(r)\|^2.
 \label{eq:tpc42-one-packet-norm}
\end{equation}
When this is summed over \(a\), uniqueness of \(\alpha\) for fixed
\((j,r)\) turns the right side into
\eqref{eq:tpc42-fixed-k-actual-diagonal}.  Expanding
\eqref{eq:tpc42-fixed-k-synthesis} gives the last identity.
\end{proof}

Thus the genuine Hilbert-valued residue variance is already exactly the
atomic diagonal.  No scalar cancellation is used, and changing any
M\"obius sign inside one residue fiber leaves
\eqref{eq:tpc42-fixed-k-packed-diagonal} unchanged.

The all-residue indexing is essential for the unconditional physical
identity.  TPC--41's scalar progression theorem is stated on
\((\Z/M\Z)^\times\), so it controls bounded postcompositions of the unit
sector only.  If an independently justified regular-packet pruning makes
the nonunit packets vanish, the starred version follows as a special
case; otherwise the nonunit packets are an additional scalar-invisible
sector, not terms that may be silently discarded.

\subsection{The complete minimal-folded identity}

For a folded bucket \(0\le s\le L_B\), define
\begin{equation}
 [\mathbf H_s^L(r)]_{\gamma,j}
 :=\sum_{\alpha:m_\alpha j=r}
 c^L_{\alpha,\gamma,j}u_{\alpha,j,s},
 \label{eq:tpc42-folded-H}
\end{equation}
and its target-residue packet
\begin{equation}
 \mathbf W_{s,a}^L
 :=\sum_{\substack{r\equiv a-h\!\!\!\pmod M}}
 \mathbf H_s^L(r).
 \label{eq:tpc42-folded-residue-packet}
\end{equation}
Although \(u_{\alpha,j,s}\) contains one or two aliases, both target
integers have residue \(m_\alpha j+h\pmod M\).  The proof of
\cref{thm:tpc42-fixed-alias-orthogonality} therefore applies without
change to \(\mathbf H_s^L(r)\).

\begin{theorem}[Exact folded residue decomposition]
\label{thm:tpc42-folded-residue-identity}
The actual minimal bank satisfies
\begin{align}
 \boxed{
 \cE_{\rm same,L}
 =\sum_{s=0}^{L_B}\sum_{a\bmod M}
   \|\mathbf W_{s,a}^L\|^2,}
 \label{eq:tpc42-folded-packed-energy}\\
 \boxed{
 \cE_{\rm fold,L}
 =\sum_{s=0}^{L_B}
   \left\|\sum_{a\bmod M}\mathbf W_{s,a}^L\right\|^2,}
 \label{eq:tpc42-folded-coherent-synthesis}\\
 \boxed{
 \cF_L
 =\sum_{s=0}^{L_B}
   \sum_{\substack{a,b\bmod M\\a\ne b}}
   \langle\mathbf W_{s,a}^L,\mathbf W_{s,b}^L\rangle.}
 \label{eq:tpc42-F-as-cross-residue}
\end{align}
The same three identities hold on the right.
\end{theorem}

\begin{proof}
For fixed \(s\), different \(r\)'s in one target residue give orthogonal
vectors.  Therefore
\begin{align}
 \sum_a\|\mathbf W_{s,a}^L\|^2
 &=\sum_r\|\mathbf H_s^L(r)\|^2
 \notag\\
 &=\sum_{\gamma,j,\alpha}
   |c^L_{\alpha,\gamma,j}|^2|u_{\alpha,j,s}|^2.
 \label{eq:tpc42-one-folded-packet-energy}
\end{align}
Summing over \(s\) proves \eqref{eq:tpc42-folded-packed-energy}.  At a
fixed coordinate \((\gamma,j)\), summing all target residues is the same
as summing all rows, which gives
\eqref{eq:tpc42-folded-coherent-synthesis}.  Finally, if
\(\alpha\ne\beta\) occur at the same \((\gamma,j)\), then
\begin{equation}
 m_\alpha j+h\not\equiv m_\beta j+h\pmod M;
 \label{eq:tpc42-distinct-rows-distinct-residues}
\end{equation}
otherwise the proof of
\cref{thm:tpc42-fixed-alias-orthogonality} would force
\(\alpha=\beta\).  Consequently the ordered cross-residue terms are
exactly the ordered cross-row terms defining \(\cF_L\), proving
\eqref{eq:tpc42-F-as-cross-residue}.
\end{proof}

Combining \eqref{eq:tpc42-folded-packed-energy} with the inherited
row-diagonal theorem gives the unconditional vector estimate
\begin{equation}
 \sum_{s,a}
 \left(
  \|\mathbf W_{s,a}^L\|^2+
  \|\mathbf W_{s,a}^R\|^2
 \right)
 \ll X^{o(1)}Q^2JK_B.
 \label{eq:tpc42-packed-vector-target}
\end{equation}
This is a diagonal-scale Hilbert-valued residue theorem for the actual
packet.  The unresolved operation is no longer packing integers inside
one residue; it is the coherent summation over different residues in
\eqref{eq:tpc42-folded-coherent-synthesis}.
